\documentclass[11pt,a4paper]{article}
\usepackage[T1]{fontenc}
\usepackage[utf8]{inputenc}
\usepackage[margin=23mm]{geometry}
\usepackage{amsmath,amssymb,mathtools,booktabs,longtable,array}
\usepackage{xurl}
\usepackage[colorlinks=true,linkcolor=blue,urlcolor=blue]{hyperref}
\usepackage{enumitem}
\setlist{itemsep=2pt,topsep=4pt}
\setlength{\parindent}{0pt}
\setlength{\parskip}{6pt}
\setlength{\emergencystretch}{3em}
\newcommand{\R}{\mathbb{R}}
\newcommand{\E}{\mathbb{E}}
\newcommand{\pos}[1]{\left[#1\right]_{+}}
\DeclareMathOperator{\CVaR}{CVaR}
\title{Real-Time MPC and Electricity-Price Forecasting\\Mathematical Formulations and Initial Experiments}
\author{PV--BESS energy-market research repository}
\date{Implementation snapshot: 4 November 2026}
\begin{document}
\maketitle
\begin{abstract}
This document describes the implemented real-time physical controllers for a
PV plant with a shared battery and grid connection, and the design of the initial
Italian electricity-price forecasting experiments. It distinguishes the shared
open-loop ensemble plan, conditional adaptive recourse, the two-stage settlement
price oracle, and ex-post realized accounting. It also states the feature and
training constructions, empirical price scenarios, evaluation metrics, and
native-resolution restrictions. Historical experiments are retrospective and
are not presented as evidence that these policies or data were available at the
time. Long-horizon controller results and final neural PV models remain pending.
\end{abstract}
\tableofcontents
\newpage

\section{Scope, events and information}
The system contains one PV plant and an AC-coupled battery behind a common point
of interconnection (POI). There is no local load. The auction model participates
in the day-ahead auction (DA) and up to three intraday auctions (IDA1--IDA3).
Each signed order is assumed to execute fully. Upon submission its quantity is
an immutable commitment; at publication its realized cashflow is booked once.
An unpublished commitment is already fixed but not yet priced. The physical
controller cannot change these submitted positions. Later auction decisions can
change because the controller has changed the actual battery state.

A \emph{continuous replay} carries the actual state, commitments and degradation
memory chronologically across delivery days. It refers to repeated auction
reoptimization and physical control, not a model of continuous order-book
trading. Forecast objectives, value-to-go terms and risk penalties are planning
quantities and are not booked as realized cash. The core auction formulation is
in the companion \path{docs/mathematical_model.tex}; here we focus on the newer
real-time models and the price inputs.

At each physical decision, the controller starts from measured or estimated
battery energy and known commitments, obtains an eligible PV forecast and,
when applicable, a settlement-price forecast, and solves over the remaining
physical horizon. It implements only the first action and repeats at the next
interval. Weather forecast cycles, lead times, their assumed release lag and
source provenance are tracked; in the matched future comparison the newest
eligible cycle covering the required interval is selected. Reviewed execution
PV can contain synthetic replacements or configured night zeros, not solely
independent meter observations. An oracle is explicitly noncausal when it uses
realized future settlement or clearing prices.

\section{Notation and shared physical constraints}
Let $t=0,\ldots,T-1$ index the intervals remaining at one decision and
$\Delta_t$ denote their duration in hours. A local day usually has 96
quarter-hour intervals, with 92 or 100 on daylight-saving transitions. Let
$s\in\mathcal S$ index PV/price scenarios with $w_s\ge0$ and
$\sum_s w_s=1$. All power is in MW, energy and commitments in MWh, and prices
in EUR/MWh. A positive grid power denotes export.

\begin{longtable}{>{\raggedright\arraybackslash}p{0.20\textwidth}>{\raggedright\arraybackslash}p{0.72\textwidth}}
\toprule Symbol & Meaning \\
\midrule\endhead
$p_{s,t}$ & Forecast available PV power in scenario $s$. \\
$c_{s,t},d_{s,t}$ & Battery charge and discharge power; their information ownership depends on the controller. \\
$u_{s,t}\in[0,1]$ & Fraction of available PV used; the remainder is curtailed. \\
$e_{s,t}$ & Battery energy at the interval boundary. \\
$g_{s,t}$ & Net power exported through the POI. \\
$q_t$ & Previously submitted net auction commitment for this interval. \\
$b^+_{s,t},b^-_{s,t}$ & Positive and negative deviations from the fixed commitment. \\
$D_s,I_s,K_s,R_s$ & Pathwise wear, import cost, terminal inventory cost and settlement revenue. \\
\bottomrule
\end{longtable}

For each path, the implementation enforces
\begin{align}
 g_{s,t} &= p_{s,t}u_{s,t}+d_{s,t}-c_{s,t},
 &-P^{\mathrm{imp}} &\le g_{s,t}\le P^{\mathrm{exp}},\\
 e_{s,t+1} &= (1-\sigma)^{\Delta_t}e_{s,t}
       +\eta_c\Delta_t c_{s,t}-\frac{\Delta_t}{\eta_d}d_{s,t},
 &E^{\min}&\le e_{s,t}\le E^{\max},\\
 0\le c_{s,t}&\le P^{\max}z_{s,t},
 &0\le d_{s,t}&\le P^{\max}(1-z_{s,t}),\quad z_{s,t}\in\{0,1\},\\
 b^+_{s,t}-b^-_{s,t}&=\Delta_t g_{s,t}-q_t,
 &b^+_{s,t},b^-_{s,t}&\ge0.
\end{align}
Initial battery energy is fixed to the current state. POI bounds, energy
bounds, optional terminal constraints and battery-mode exclusivity apply on
every scenario. $\sigma$ is hourly self-discharge and $\eta_c,\eta_d$ are AC
charge/discharge efficiencies. Any configurable import charge is an epigraph
of $\max(-g_{s,t},0)$; the project does not impose an extra grid-usage fee
merely because the battery imports.

With a throughput wear coefficient $\gamma$, one supported wear model is
$D_s=\gamma\sum_t\Delta_t(c_{s,t}+d_{s,t})$. The alternative segmented Xu
cycle-depth implementation includes its own segment-ordering and carried cycle
memory; $D_s$ below denotes whichever configured degradation model is used.
See \path{src/optimization/degradation.py} and the companion mathematical
model for its detailed constraints. A soft terminal target is, for example,
$K_s=\rho\lvert e_{s,T}-e^{\mathrm{target}}\rvert$; linear or configured
piecewise-linear continuation values are alternative terminal treatments.

\section{Common-control deterministic and ensemble MPC}
\subsection{Control ownership}
The deterministic controller uses one forecast path. The ordinary ensemble MPC
uses multiple PV paths but a \emph{single open-loop physical schedule}:
\begin{equation}
(c_{s,t},d_{s,t},u_{s,t},z_{s,t},e_{s,t})
 =(c_t,d_t,u_t,z_t,e_t)\quad\text{for all }s.
\end{equation}
Grid injections and deviations still differ with $p_{s,t}$. All scenario POI
constraints must hold. In a closed-loop replay the common schedule is not
committed wholesale: only its first action is executed, and a new schedule is
computed after the next observation.

For the price-blind penalty controller, define scenario-weighted absolute
imbalance
\begin{equation}
 A=\sum_s w_s\sum_t(b^+_{s,t}+b^-_{s,t}).
\end{equation}
The solver minimizes
\begin{equation}\label{eq:blind}
 \kappa_{\mathrm{blind}} A + D + I + K,
\end{equation}
where $D$ and $K$ are shared trajectory costs and $I$ is expected import cost.
The penalty case in the research configuration uses $\kappa_{\mathrm{blind}}=500$
EUR/MWh. This is a trade-off between expected imbalance and other costs; it is
not a lexicographic requirement that $A$ attain its minimum. With a positive
absolute-deviation penalty the split is tight at optimum, subject to solver
numerical tolerances. Equal scenario/interval PV values may share identical
interval-support variables; this exact deduplication does not delete members or
relax the POI envelope.

\subsection{Two-stage realized-settlement-price oracle}
The optional oracle uses actual settlement curves, while its first stage remains
price-blind. The stages are two optimizations within one MPC decision, not the
DA and intraday auctions:
\begin{enumerate}
\item Minimize $A$ subject to the shared physical constraints. Let $U$ be the
independently reconstructed feasible imbalance and $L$ the solver lower bound.
\item If stage one meets the configured optimality criterion, or the explicitly
configured time-limit fallback has a qualifying reconstructed relative gap,
maximize settlement revenue less wear, import and terminal costs:
\begin{align}
 \max\quad &R-D-I-K,\\
 \text{subject to}\quad&A\le U+\beta U+\delta,\\
 R={}&\sum_s w_s\sum_t\bigl(\pi_t^+b^+_{s,t}-\pi_t^-b^-_{s,t}\bigr).
\end{align}
Here $\beta$ and $\delta$ are configured relative and absolute extra
imbalance budgets. The second stage is restricted by the feasible stage-one
incumbent, which need not be a proven minimum when the fallback is used.
\end{enumerate}
For positive $U$, the reconstructed gap is $(U-L)/U$; the implementation also
handles near-zero objectives explicitly. When enabled in selected studies, the
fallback accepts a time-limited incumbent only if its gap is at most 1\%.
The standard relative solver MIP gap is $10^{-4}$ and no absolute MIP gap is
configured. These studies allocate a separate 10-second optimization budget to
each stage; building and validating models add wall time. A failed or invalid
second-stage refinement retains the validated first-stage plan. A qualified
incumbent is not described as a certified optimum. Direction-specific settlement
prices require a binary sign restriction to prevent simultaneous artificial
positive and negative deviations. The equal-price case uses signed settlement.

\section{Forecast-economic and adaptive ensemble MPC}
\subsection{Scenario-based forecast economics}
The research economic controllers use forecast, rather than realized,
settlement prices $\hat\pi^+_{s,t}$ and $\hat\pi^-_{s,t}$. Let
\begin{equation}
 R_s=\sum_t(\hat\pi^+_{s,t}b^+_{s,t}-\hat\pi^-_{s,t}b^-_{s,t}),
 \qquad L_s=D_s+I_s-R_s.
\end{equation}
The implemented objective is
\begin{equation}\label{eq:economic}
 \min\sum_s w_s\left(D_s+I_s+K_s-R_s+
 \kappa_{\mathrm{econ}}\sum_t(b^+_{s,t}+b^-_{s,t})\right)
 +\lambda\CVaR_\alpha(L).
\end{equation}
The economic imbalance coefficient is separate from the price-blind penalty
and defaults to zero in these research cases. $L_s$ excludes the artificial
imbalance penalty and terminal salvage, and does not represent uncertain
receipts on unpriced auction commitments. At the configured risk-neutral
baseline $\lambda=0$. When enabled, weighted monetary CVaR has the linear form
\begin{equation}
 \CVaR_\alpha(L)=\min_{\eta\in\R,\,\xi_s\ge0}
 \left\{\eta+\frac{1}{1-\alpha}\sum_s w_s\xi_s:
 \xi_s\ge L_s-\eta\right\}.
\end{equation}
The common-control economic ensemble controller uses the same physical
ownership as the ordinary ensemble MPC. A deterministic economic controller
has one PV path. Forecast settlement paths and PV members are combined under
the declared scenario construction; the initial workflow does not learn a
joint multi-auction/PV/settlement distribution.

\subsection{Adaptive PV-history recourse}
The adaptive controller optimizes objective~\eqref{eq:economic} while allowing
future controls to depend on disclosed PV histories. Let $n(s,t)$ be the node
assigned to scenario $s$ at time $t$. If two scenarios share a node then their
charge, discharge, PV utilization and battery mode are identical:
\begin{equation}
 n(s,t)=n(s',t)\Longrightarrow
 (c_{s,t},d_{s,t},u_{s,t},z_{s,t})
 =(c_{s',t},d_{s',t},u_{s',t},z_{s',t}).
\end{equation}
All scenarios share the first node and the current command. Nodes are formed
from quantized observed PV prefixes up to a configured revelation lag; branch
points occur only at configured future indices. They do not use future PV or a
scenario label. Battery state follows each node's control history, and nodes do
not recombine distinct prior battery histories. The current implementation does
not reveal simulated settlement prices along a recourse path.

The prepared long study uses branching every four quarter-hour intervals, a
one-interval observation lag and 0.1 MW PV-history bins; these are study
settings, not structural constants. This is a restricted information policy,
not perfect foresight. A retained representative path is an auction shadow
projection and is not a prediction of which scenario will occur. Only the
common first action is executed. Adaptive and ordinary ensemble models should
be compared with matched objective, scenarios, information, state and physical
settings before attributing a difference to recourse itself.

\section{Execution, settlement and information comparisons}
The grid-tracking baseline follows the committed grid position using its
configured interval PV feedback and physical protection. MPC uses its own
configured current-PV input, often the forecast. These information sets differ.
Physical clipping or safety overrides and solver fallbacks are recorded.

For a realized interval, $\hat g_t$ is the executed grid power and
$\hat b_t=\Delta_t\hat g_t-q_t$. Actual settlement is calculated from realized
positive and negative deviations and realized settlement prices; actual auction
cash is booked only on publication. Realized operating profit combines those
cashflows and configured fees, imports and wear. It is not the optimizer's
forecast objective. Since each policy carries its own state and can submit
later auction orders, a profit comparison reflects the whole closed-loop policy.
Ending SOC and terminal treatment must be reported alongside profit.

\section{Initial price forecasting experiments}
\subsection{Data and feature construction}
The initial retrospective experiments forecast the NORD DA, IDA1--IDA3 and
imbalance-settlement prices. DA training starts in 2024 in the specified ridge
study; its IDA and settlement training starts in 2025. A separate historical
intraday experiment also used 2024 hourly exports as coarse labels expanded
onto quarter-hours. Those repeated rows are not independent quarter-hour
products. The newer search selects native hourly and native quarter-hourly
products \emph{before} any expansion and fits separate models for each market
and resolution. Native price products are held constant within their duration,
not linearly interpolated.

The initial price-only feature vector contains historical price lags and local
calendar terms: sine/cosine representations of hour, weekday and annual phase,
plus a weekend indicator. The European-context ridge study uses 24, 48, 168 and
336 hour lags, with configured same-clock fallback when a lag is unavailable.
Its NORD fundamentals add zonal load, total-generation, solar and onshore-wind
forecasts. Additional ablations add PSV gas, Milan/Turin temperature forecasts,
French nuclear and hydro generation and France--Italy commercial exchanges, or
Swiss aggregate balance. A combined set includes these groups. An August-only
set includes Swiss--Italy commercial schedules, directional \emph{intraday}
NTC and the Italian-NORD/Swiss spot spread. These are not physical border-flow
measurements. No individual farm PV ensemble is a zonal-price feature.

The initial context run applies stated hypothetical publication policies:
French/Swiss realized observations and PSV are shifted 48 hours; Swissgrid
border observations use a nine-day shift with an assumed weekly release lag.
A separate retrospective experiment uses previous-day realized information and
a deliberately permissive source-day publication placeholder. Exact release
clocks and revisions are not established by the download timestamp. Weather
hindcasts are allowed in the stated retrospective temperature comparison.
Missing source observations are retained as gaps, not converted to zero.

\subsection{Ridge model, temporal split and residual paths}
For training rows $i$, let $x_i$ be the selected lag, calendar and optional
fundamental features, $y_i$ the observed price and $\tilde x_i$ the features
standardized by training-set mean and scale. The fitted point model is
\begin{equation}
 \hat y_i=\beta_0+\tilde x_i^{\mathsf T}\beta,\qquad
 (\hat\beta_0,\hat\beta)=\arg\min_{\beta_0,\beta}
 \sum_{i\in\mathcal T_{\mathrm{train}}}
 (y_i-\beta_0-\tilde x_i^{\mathsf T}\beta)^2+\alpha\lVert\beta\rVert_2^2.
\end{equation}
The intercept is unpenalized. The initial European-context comparison sets
$\alpha=20$ and uses a 14-day chronological calibration tail. A predictor
trained before the calibration period generates residuals on that period; the
deployed point predictor is then refitted on all eligible labels through the
fold cutoff. The calibration residuals do not come from the refitted model's
in-sample predictions. Training labels and the feature-vintage eligibility rules
are recorded separately from evaluation periods.

Whole complete local-day residual blocks $\varepsilon_{d,t}$ provide empirical
price paths. For a selected calibration day $d_j$, a forecast sample is
\begin{equation}
 \hat\pi^{(j)}_t=\hat y_t+\varepsilon_{d_j,t}.
\end{equation}
Sampling a block preserves its within-day residual shape, rather than drawing
independent noise at each interval. The configured DST policy controls which
day lengths or local-clock slots can be used. Residual draws for separate days
are independent; this does not model cross-day dependence or establish a joint
PV--price distribution. The separate point forecast $\hat y_t$ need not equal
the finite residual-scenario mean.

\subsection{Scores and experimental design}
For observed prices $y_i$ and point forecasts $\hat y_i$, the reported absolute
and relative errors include
\begin{align}
 \mathrm{MAE}&=\frac1N\sum_{i=1}^N|y_i-\hat y_i|,\\
 \mathrm{WAPE}(\%)&=100\frac{\sum_i|y_i-\hat y_i|}{\sum_i|y_i|},\\
 \mathrm{rMAE}&=\frac{\sum_i|y_i-\hat y_i|}{\sum_i|y_i-\hat y_i^{\mathrm{naive}}|}.
\end{align}
WAPE and rMAE are undefined when their respective denominators are zero. WAPE
is a pooled normalization, not the mean of individual percentage errors; this
avoids dividing separately by zero or near-zero prices. In the initial ridge
report, MAE is averaged across complete daily scores, while WAPE is pooled over
the same common-day intervals; these need not weight 23/25-hour days identically.
For forecast samples $x_1,\ldots,x_M$, the empirical CRPS at observation $y$
is
\begin{equation}
 \mathrm{CRPS}=\frac1M\sum_{j=1}^{M}|x_j-y|
 -\frac{1}{2M^2}\sum_{j,k=1}^{M}|x_j-x_k|.
\end{equation}
CRPS evaluates the predictive distribution; it is not the ridge training loss
or the eventual trading profit. The initial context ablation evaluates January
and August 2026 and compares declared models on common complete delivery days.
Feature groups can have different training coverage, particularly Swiss series
and the 2026 border file; therefore shared test dates alone do not guarantee
identical training samples. The results are development comparisons, not a
full-year or universally causal test.

The expanded CPU search adds ridge, Lasso, elastic net, feed-forward neural
networks, random forests, histogram gradient boosting, LightGBM and CatBoost
candidates. These are point-model candidates on the library's interval features;
the Lasso is not a reproduction of the full-curve LEAR benchmark. The first
search validates separately by native resolution and reserves later dates for
selection-fixed evaluation. Its quarter-hour later period includes August 2026,
which was inspected in an earlier ridge study, so it is not globally untouched.
Intraday and settlement searches are separate market tasks. The initial search
is not an exhaustive hyperparameter search, nor are the selected candidates yet
portable MPC forecast artifacts. The design draws on the benchmarking and
metric discussion of Lago et al.\ \cite{lago2021}; model performance and
forecast quality remain distinct from closed-loop trading value.

\section{Prepared longer controller studies and limitations}
The prepared annual study defines one continuous 2026 trajectory per farm and
case; the ensemble study defines separate full January and August trajectories.
The latter does not simulate the intervening months and therefore starts each
month from its configured state. Both define matched pvlib and eventual neural
forecast cases, with continuous commitments and latest-eligible forecast cycles.
The present reviewed PV/measurement inputs end in August 2026, cached market
series do not yet cover the whole year, and the available August ensemble lacks
the preceding 31 July cycle needed by the first DA decision. The final neural
models are being trained. Existing file presence is not proof of temporal
coverage: complete replay preflight must check each required date before a run.
No full-year or full-month controller result is claimed here.

\section{Implementation references}
\begin{itemize}
\item Common-control penalty and oracle: \path{src/optimization/mpc.py}; shared battery constraints: \path{src/optimization/physical.py}.
\item Forecast-economic and adaptive model: \path{src/optimization/adaptive_mpc.py}; information nodes: \path{src/scenario_research.py}.
\item Price features and fitting: \path{src/forecasting/energy_prices.py}, \path{src/forecasting/energy_price_features.py}, and \path{scripts/fit_energy_price_forecast.py}.
\item Native-resolution model search: \path{scripts/search_energy_price_models.py}, \path{src/forecasting/price_search.py}, and \path{src/forecasting/price_search_models.py}.
\item Experiment configurations and current input limitations: \path{configs/research/price_european_context.yaml} and \path{docs/reviews/long_horizon_2026_study.md}.
\end{itemize}
\begin{thebibliography}{9}
\bibitem{lago2021} J. Lago, G. Marcjasz, B. De Schutter and R. Weron,
``Forecasting day-ahead electricity prices: A review of state-of-the-art algorithms,
best practices and an open-access benchmark,'' \emph{Applied Energy}, 293,
116983, 2021. \url{https://doi.org/10.1016/j.apenergy.2021.116983}.
\end{thebibliography}
\end{document}
